% CCSS 7.SP.B.4 — Comparing Populations with Measures of Center and Variability
\section{Comparing Populations with Measures of Center and Variability}

\begin{learningGoals}
\begin{itemize}
    \item Use mean, median, range, IQR, and MAD to compare populations
    \item Support conclusions with data, graphs, and reasoning
\end{itemize}
\end{learningGoals}

\begin{conceptBox}{Measures of Center and Variability}
\begin{itemize}[leftmargin=*]
    \item \textbf{Mean}: Average
    \item \textbf{Median}: Middle value
    \item \textbf{Range}: Difference between highest and lowest
    \item \textbf{Interquartile Range (IQR)}: Range of the middle half
    \item \textbf{Mean Absolute Deviation (MAD)}: Average distance from mean
\end{itemize}
\end{conceptBox}

\begin{workedExample}{Comparing Two Groups}
Group A: mean $60$, median $62$, range $20$, MAD $5$
Group B: mean $65$, median $66$, range $15$, MAD $4$

Group B is higher on average and less variable.
\answerHighlight{Group B has higher center and less spread.}
\end{workedExample}

\tipBox{Use multiple measures to compare groups, not just one.}

\begin{practiceBox}{Comparing Populations}
\resetProblems

\prob Group A: mean $70$, Group B: mean $75$. Which is higher?
\answerExplain{Group B}{Mean $75$ is higher than $70$.}

\prob If Group A has range $30$, Group B has range $20$, which is more variable?
\answerExplain{Group A}{A larger range means more variability.}

\prob What does a higher MAD mean?
\answerExplain{More spread out data}{Higher MAD means data points are farther from the mean.}

\prob If two groups have the same median but different means, what does this mean?
\answerExplain{Means measure average, medians measure middle}{Groups may have similar middle values but different averages.}

\prob Why use IQR to compare groups?
\answerExplain{Shows spread of middle half}{IQR focuses on the central part of the data, less affected by outliers.}

\end{practiceBox}
